%% GRT-Parametrized Nine Faces of the Collision Residue.
%% Constructs the torsor realization of Face_K(A) over GRT_1(K), promotes
%% the historical ``seven faces'' to rational finite-gauge representatives
%% where verified, and adjoins Brown's motivic-period face F8 and
%% Willwacher's operadic face F9 after the relevant scalar extension.
%%
%% Vol II theory sequence. Depends on factorization_swiss_cheese,
%% sc_chtop_heptagon, and Vol I seven_faces (standalone).
%%
%% Macro safety: all control sequences used in this file are guarded
%% by \providecommand here; no \newcommand.

\providecommand{\SCchtop}{\mathsf{SC}^{\mathrm{ch},\mathrm{top}}}
\providecommand{\FM}{\mathrm{FM}}
\providecommand{\Conf}{\mathrm{Conf}}
\providecommand{\Ran}{\mathrm{Ran}}
\providecommand{\MC}{\mathrm{MC}}
\providecommand{\End}{\mathrm{End}}
\providecommand{\cA}{\mathcal{A}}
\providecommand{\cB}{\mathcal{B}}
\providecommand{\cF}{\mathcal{F}}
\providecommand{\cL}{\mathcal{L}}
\providecommand{\cW}{\mathcal{W}}
\providecommand{\Walg}{\mathcal{W}}
\providecommand{\fg}{\mathfrak{g}}
\providecommand{\fgrt}{\mathfrak{grt}_1}
\providecommand{\GRT}{\mathrm{GRT}_1}
\providecommand{\GT}{\mathrm{GT}_1}
\providecommand{\Lie}{\mathsf{Lie}}
\providecommand{\Ass}{\mathsf{Ass}}
\providecommand{\Face}{\mathrm{Face}}
\providecommand{\Fsuper}{\mathrm{Face}^{\mathrm{super}}}
\providecommand{\Gsuper}{\mathrm{GRT}^{\mathrm{super}}}
\providecommand{\Res}{\operatorname{Res}}
\providecommand{\coll}{\mathrm{coll}}
\providecommand{\Vir}{\mathrm{Vir}}
\providecommand{\id}{\mathrm{id}}
\providecommand{\GC}{\mathsf{GC}}
\providecommand{\Def}{\mathrm{Def}}
\providecommand{\Tam}{\mathrm{Tam}}
\providecommand{\KZ}{\mathrm{KZ}}
\providecommand{\rat}{\mathrm{rat}}
\providecommand{\mot}{\mathrm{mot}}
\providecommand{\MZV}{\mathrm{MZV}}
\providecommand{\grt}{\mathrm{grt}}
\providecommand{\DNP}{\mathrm{DNP}}
\providecommand{\GZ}{\mathrm{GZ}}
\providecommand{\FFR}{\mathrm{FFR}}
\providecommand{\STS}{\mathrm{STS}}
\providecommand{\Gaud}{\mathrm{Gaud}}
\providecommand{\Sug}{\mathrm{Sug}}
\providecommand{\Bbbk}{\mathbb{k}}

\chapter{The GRT-Parametrized Nine Faces of the Collision Residue}
\label{ch:grt-seven-faces}

%%%%%%%%%%%%%%%%%%%%%%%%%%%%%%%%%%%%%%%%%%%%%%%%%%%%%%%%%%%%%%%%%%%%%%%%%%%%%
\section*{Prologue}
%%%%%%%%%%%%%%%%%%%%%%%%%%%%%%%%%%%%%%%%%%%%%%%%%%%%%%%%%%%%%%%%%%%%%%%%%%%%%

Seven papers named the same object. The object is a single element
of the $E_1$ ordered convolution algebra,
\[
 r(z) \;=\; \Res^{\coll}_{0,2}\bigl(\Theta_{\cA}\bigr)
 \;\in\; \cA^{!}\otimes\cA^{!}[\![z^{-1}]\!];
\]
seven literatures attach seven different normalisations, seven
different Casimir bases, seven different connection conventions.
Vol~I treated the seven names as seven vertices of a star and proved
a seven-way equivalence by routing every edge through the bar hub.
The present chapter rewrites the star as a \emph{torsor}. The
secret party underneath the star is the Drinfeld
Grothendieck--Teichm\"uller group $\GRT(K)$ acting by
associator transport. Each ``face'' is an associator-coordinate
chart over a coefficient field $K$; the seven rational finite-gauge
representatives have two canonical
companions, the Brown motivic face $F_8$ and the Willwacher
operadic face $F_9$, whose comparison is meaningful only after scalar
extension to a common coefficient field. The pentagon is a torsor, the
torsor is a scheme over $\GRT$, and ``seven'' is the historical
rational finite-gauge slice.

The reorganisation closes four persistent confusions. The
landscape-census chained equality
$k\cdot\Omega_{\mathrm{tr}}/z = \Omega/((k+h^{\vee})z)$ is not a
rational-function identity; it is a Casimir rescaling
belonging to the finite inner gauge subgroup
$\GRT^{\mathrm{fin}}\subset\GRT$, a tensor identity at different
normalisations. The F1$\leftrightarrow$F4 identification is an
injection of $r(z)$ into Drinfeld associator data; the MZV tail
$\Phi_{\KZ}-\Phi_{\rat}$ is invisible to $r(z)$. The F7
entry is ambiguous: F7 (Gaudin simple-pole) and F7$'$ (class-M top
$A_\infty$ $m_3$) coincide on class L but separate on class M.
The Heisenberg-vs-symplectic-fermion asymmetry is a
$\mathbb{Z}/2$-factor on collision residues, absorbed into the
super-variant $\Gsuper_K = \GRT(K)\rtimes(\mathbb{Z}/2)^{|\mathrm{odd}|}$.

The chapter is organised as follows. Section~\ref{sec:grt-setup}
defines $\Face_K(\cA)$ as the scheme of chain-level presentations of
$r(z)$ compatible with the bar-intrinsic dGLA, and proves that it
carries a canonical $\GRT$-torsor structure.
Section~\ref{sec:grt-nine} enumerates the nine faces F1--F9.
Section~\ref{sec:grt-torsor-thm} states the master theorem:
$\Face_K(\cA)$ is a $\GRT(K)$-torsor after a coefficient field and
formality base point are fixed. Section~\ref{sec:f1-f4}
proves the F1$\leftrightarrow$F4 injection and the invisibility of
MZV higher weights. Section~\ref{sec:f7-disambig} disambiguates F7
and F7$'$ at class L vs class M. Section~\ref{sec:grtfin}
reclassifies the Casimir rescaling as a $\GRT^{\mathrm{fin}}$ gauge.
Section~\ref{sec:super} installs the super-variant.
Sections~\ref{sec:f8-motivic} and~\ref{sec:f9-willwacher}
install $F_8$ (Brown motivic) and $F_9$ (Willwacher operadic)
as higher cochain faces. Section~\ref{sec:f8-f9-dual} records the
conditional comparison identifying $F_8$ and $F_9$ as geometric and
operadic incarnations of one orbit after scalar extension.

\begin{convention}[Coefficient field and binary residue]
\label{conv:grt-coefficient-field}
\index{GRT torsor!coefficient field}%
\index{collision residue!binary degree}%
All torsor statements are made over a characteristic-zero coefficient
field $K$ after choosing a formality base point. The notation
$\Face_K(\cA)$ means the chain-level face space after scalar extension
to $K$, and the acting group is $\GRT(K)$. The historical faces
$F_1,\ldots,F_7$ are rational finite-gauge presentations only where the
transport is explicitly binary or connection-level. Brown's face
$F_8$ lives over the motivic period algebra, or over its real period
realisation; Willwacher's face $F_9$ is defined over $\mathbb Q$ at the
graph-complex Lie-algebra level. The comparison $F_8\leftrightarrow
F_9$ is therefore a comparison after scalar extension to a common
coefficient field, not a statement that $F_8$ is a rational point of
$\GRT(\mathbb Q)$.

The symbol $r(z)$ is reserved for the binary degree-two collision
residue. MZV and graph-complex data occur in higher cochain
operations, beginning with the arity-three/twisting component
$\tau_{(3)}$ or the corresponding $A_\infty$ operation $m_3$, not in
the cohomological binary residue itself.
\end{convention}

%%%%%%%%%%%%%%%%%%%%%%%%%%%%%%%%%%%%%%%%%%%%%%%%%%%%%%%%%%%%%%%%%%%%%%%%%%%%%
\section{Setup: \texorpdfstring{$\Face_K(\cA)$}{Face_K(A)} as a GRT-torsor}
\label{sec:grt-setup}
%%%%%%%%%%%%%%%%%%%%%%%%%%%%%%%%%%%%%%%%%%%%%%%%%%%%%%%%%%%%%%%%%%%%%%%%%%%%%

Fix a chirally Koszul chiral algebra $\cA$ in the standard landscape
and a characteristic-zero coefficient field $K$ as in
Convention~\ref{conv:grt-coefficient-field}.
Let $\Theta_{\cA}^{E_1}\in\mathfrak{g}^{E_1}_{\cA}$ denote the
universal Maurer--Cartan element in the ordered convolution
algebra (Vol I, \S\text{shadow-tower}); its degree-$2$ collision
residue is
\[
 r(z) \;=\; \Res^{\coll}_{0,2}(\Theta_{\cA}^{E_1}).
\]

\begin{definition}[\label{def:face-space}Face space]
A \emph{chain-level presentation} of $r(z)$ compatible with the
bar-intrinsic dGLA is the datum of
\begin{enumerate}[(a)]
\item a chain model $\mathfrak{m}$ of the $E_1$-chiral convolution
dGLA, together with a chain-level realisation of $\Theta_{\cA}^{E_1}$
as a strict $\MC$-element in $\mathfrak{m}$;
\item a Casimir basis $(C_\alpha)_{\alpha\in I}$ for the
binary-degree component of $\mathfrak{m}$, inducing a coordinate
presentation $r(z) = \sum_\alpha r^\alpha(z)\,C_\alpha$;
\item a connection convention $\nabla$ with respect to which
$\Theta_{\cA}^{E_1}$ satisfies the flat-section equation
$\nabla\Phi = 0$;
\item a Drinfeld associator $\Phi\in\GRT(K)$ providing the
isomorphism of $\infty$-operads
$\Tam_\Phi\colon B(\cA)\xrightarrow{\sim} B_\Phi(\cA)$ (Tamarkin).
\end{enumerate}
The set of equivalence classes of such presentations after scalar
extension to $K$, under strict $\MC$-gauge equivalence in
$\mathfrak{m}$, is the \emph{face space} $\Face_K(\cA)$.
\end{definition}

\begin{remark}[Distinguished base point; $F_1$ is identity]
\label{rem:f1-base-point}
The bar hub $F_1$, i.e. the twisting morphism
$\pi_{\cA}\colon B(\cA)\to\cA$ with deconcatenation coproduct and
identity Drinfeld datum $\Phi = \id$, is a distinguished point of
$\Face_K(\cA)$. Fixing $F_1$ as origin converts the torsor into a
group; everywhere in the chapter we use $F_1$ as the identity coset.
\end{remark}

\begin{proposition}[\label{prop:grt-action-face}GRT-action on $\Face_K(\cA)$;
\ClaimStatusConditional]
After fixing $K$, a formality base point, and a chain model, the group
$\GRT(K)$ acts freely and transitively on the associator-coordinate
component of $\Face_K(\cA)$. The abstract free-transitive action is
Tamarkin--Drinfeld--Willwacher formality; the concrete assertion that
every named face is a $K$-point is a per-face statement with the
coefficient-field scope of
Convention~\ref{conv:grt-coefficient-field}.
\end{proposition}

\begin{proof}
Tamarkin's theorem (Tamarkin 1998; Gerstenhaber--Voronov;
reformulated in Willwacher 2015, Thm.~1.1) constructs a free and
transitive action of $\GRT(K)$ on the space of
quasi-isomorphism classes of formality transfers
$E_2\xrightarrow{\sim} B(\Ass^{!})$, which in turn acts on the
bar-intrinsic dGLA of any associative chiral algebra through the
canonical functor $\Ass\to\ChirAss$. Explicitly, for
$\Phi\in\GRT(K)$ the action on a presentation
$(\mathfrak{m},r,\nabla,\Phi_0)$ returns the presentation with
associator $\Phi\cdot\Phi_0$, Casimir basis pulled back by the
finite inner gauge of $\Phi$, and connection convention
$\nabla^\Phi = \Phi\cdot\nabla$. Freeness is the statement that
$\GRT(K)$ acts without fixed points on formality transfers, which is
Drinfeld's 1991 Leningrad theorem. Transitivity is the pentagon axiom
combined with Tamarkin--Kontsevich formality on the
associator-coordinate component. The final step from an abstract
associator coordinate to a named face is not uniform in $K$: F8 and F9
require the scalar extensions specified in
Convention~\ref{conv:grt-coefficient-field}.
\end{proof}

\begin{corollary}[\label{cor:face-torsor}Torsor structure]
Once $F_1$ is fixed as base point, the map $\Phi\mapsto\Phi\cdot F_1$
is a bijection from $\GRT(K)$ to the associator-coordinate component of
$\Face_K(\cA)$. That component is a principal homogeneous
$\GRT(K)$-space, canonically a torsor.
\end{corollary}

\begin{remark}[Why a torsor and not a star]
\label{rem:torsor-vs-star}
Vol I's seven-face programme described the set of seven
presentations as a star with the bar hub at the centre, routing
every pairwise identification through $F_1$. The star picture is
correct as an \emph{accessibility diagram}; it does not record the
groupoid structure of the gauge transports between non-hub
vertices. The torsor picture records all gauge transports
simultaneously: the edge $F_i\to F_j$ is an element
$\Phi_{ji}\in\GRT(K)$ after scalar extension to a common coefficient
field, and the star is the subgraph of
edges through $F_1 = \id$.
\end{remark}

\begin{remark}[First-principles alignment]
The ghost theorem of the old star statement is the seven-way
equivalence of presentations; the precise error is the suppression
of the $\GRT$-action that realises those equivalences; the correct
relationship is torsor action, with the hub playing only the role
of a distinguished base point.
\end{remark}

%%%%%%%%%%%%%%%%%%%%%%%%%%%%%%%%%%%%%%%%%%%%%%%%%%%%%%%%%%%%%%%%%%%%%%%%%%%%%
\section{Nine-face enumeration}
\label{sec:grt-nine}
%%%%%%%%%%%%%%%%%%%%%%%%%%%%%%%%%%%%%%%%%%%%%%%%%%%%%%%%%%%%%%%%%%%%%%%%%%%%%

The nine face coordinates are listed in Table~\ref{tab:nine-faces}.
Seven are historical rational/finite-gauge coordinates where verified;
F8 and F9 enter as higher cochain coordinates after the scalar
extensions of Convention~\ref{conv:grt-coefficient-field}.

\begin{table}[h]
\centering
\begin{tabular}{@{}lll@{}}
\toprule
Face & Name & Reference \\
\midrule
$F_1$ & Bar hub; $\pi_{\cA}\colon B(\cA)\to\cA$, deconcatenation coproduct
 & Ginzburg--Kapranov 1994; Loday--Vallette Ch.~11 \\
$F_2$ & DNP R-twist on open colour
 & Dimofte--Niu--Py arXiv:2508.11749 \\
$F_3$ & KZ classical PVA $r$-matrix
 & De Sole--Kac arXiv:0712.1045 \\
$F_4$ & GZ26 commuting differentials
 & Gaiotto--Zeng 2026, sphere Hamiltonians \\
$F_5$ & Yangian classical $r$-matrix
 & Drinfeld 1985 (Leningrad); RTT presentation \\
$F_6$ & Gaudin Hamiltonian simple-pole
 & Feigin--Frenkel--Reshetikhin 1994 \\
$F_7$ & Sklyanin tensor; $F_7'$ class-M top $A_\infty$ operation
 & STS 1983; Mizoguchi--Ohtsuki (see \S\ref{sec:f7-disambig}) \\
$F_8$ & Brown motivic-period higher cochain face
 & Brown 2012 arXiv:1102.1312 \\
$F_9$ & Willwacher operadic face; $\fgrt$-action on $\GC_2$
 & Willwacher 2015 arXiv:1009.1654 \\
\bottomrule
\end{tabular}
\caption{Nine canonical face coordinates. Seven are the historical
rational finite-gauge presentations where verified; the last two are
motivic-period and graph-complex higher cochain coordinates.}
\label{tab:nine-faces}
\end{table}

\begin{remark}[Why ``seven'' is the wrong cardinality]
The seven-face enumeration selects exactly the coordinates that can be
written without motivic or graph-complex machinery. Once the torsor is
made visible, Brown's motivic face $F_8$ and Willwacher's operadic face
$F_9$ appear as higher cochain coordinates at the generators of
$\fgrt = \mathrm{Lie}(\GRT)$. Cutting them off is harmless for the
binary residue $r(z)$ but hides the associator-dependent tower.
\end{remark}

%%%%%%%%%%%%%%%%%%%%%%%%%%%%%%%%%%%%%%%%%%%%%%%%%%%%%%%%%%%%%%%%%%%%%%%%%%%%%
\section{The master theorem: \texorpdfstring{$\Face_K(\cA)$}{Face_K(A)} is a
$\GRT(K)$-torsor}
\label{sec:grt-torsor-thm}
%%%%%%%%%%%%%%%%%%%%%%%%%%%%%%%%%%%%%%%%%%%%%%%%%%%%%%%%%%%%%%%%%%%%%%%%%%%%%

\begin{theorem}[\label{thm:seven-faces-as-GRT-torsor}Seven faces as a
$\GRT$-torsor; \ClaimStatusConditional]
Let $\cA$ be chirally Koszul in the standard landscape and
$r(z) = \Res^{\coll}_{0,2}(\Theta_{\cA}^{E_1})$. After fixing the
coefficient field $K$ and a formality base point, the
associator-coordinate component of $\Face_K(\cA)$ is a principal
homogeneous $\GRT(K)$-space. For every pair of faces $F_i,F_j$ defined
over a common coefficient field $K$ there exists a unique associator
$\Phi_{ij}\in\GRT(K)$ such that
\[
 F_i \;=\; \Phi_{ij}\cdot F_j.
\]
The binary historical faces $F_1,\dots,F_6$ and the Gaudin/Sklyanin
simple-pole face $F_7$ are related by rational finite-gauge transports
where the Casimir and connection conventions have been fixed.
The class-M face $F_7'$, Brown's $F_8$, and Willwacher's $F_9$ are
higher cochain faces; they enter the same torsor only after the scalar
extensions named in Convention~\ref{conv:grt-coefficient-field}.
\end{theorem}

\begin{proof}
Combine Proposition~\ref{prop:grt-action-face} and
Corollary~\ref{cor:face-torsor}. The rational finite-gauge claim for
the binary historical faces amounts to showing that each such
presentation is induced from $F_1$ by an explicit finite gauge
transport after the relevant Casimir and connection convention is
fixed. The transports are:

\smallskip\noindent
\emph{$\Phi_{12}$ (Bar~$\to$ DNP R-twist).} The DNP $R$-twist
acts on the separately supplied Drinfeld spectral coproduct,
\[
\Delta_z^R = R(z)\Delta_z R(z)^{-1},\qquad
R(z) = 1 + \hbar r(z) + O(\hbar^2).
\]
The bar deconcatenation coproduct remains the structural
topological-factorisation coproduct. At order $\hbar$ the twist is
the identity on the associator, so $\Phi_{12}\in\GRT^{\mathrm{fin}}$.

\smallskip\noindent
\emph{$\Phi_{13}$ (Bar~$\to$ KZ PVA).} The KZ convention writes
$r(z) = \Omega/((k+h^\vee)z)$ in the Killing--Casimir basis; the
transport is a Casimir rescaling combined with a simple-pole shift
$k\mapsto k+h^\vee$. Both belong to $\GRT^{\mathrm{fin}}$
(Section~\ref{sec:grtfin}).

\smallskip\noindent
\emph{$\Phi_{14}$ (Bar~$\to$ GZ26).} The Gaiotto--Zeng commuting
differentials are generated by the Laplace transform of the Fay-type
kernel in the $r$-matrix presentation; the Laplace transform is the
identity on associator. The commuting property is the flatness of
$\nabla^\Phi$ on the sphere (Section~\ref{sec:f1-f4}).

\smallskip\noindent
\emph{$\Phi_{15}$ (Bar~$\to$ Yangian).} The classical limit of the
Yangian $R$-matrix lives in $\GRT^{\mathrm{fin}}$ because
$R_{\mathrm{Yang}}(u)-1 = \frac{\hbar P}{u}+O(\hbar^2)$ is a
rational function in $u$ with no higher MZV data at leading order.

\smallskip\noindent
\emph{$\Phi_{16}$ (Bar~$\to$ Gaudin).} The Gaudin residue
$H_i^{\Gaud} = \sum_{j\neq i}\Omega_{ij}/(z_i-z_j)$ is a simple-pole
reduction of $F_1$'s binary-degree component; $\Phi_{16}$ is the
finite inner gauge realising the simple-pole truncation.

\smallskip\noindent
\emph{$\Phi_{17}$ (Bar~$\to$ STS; $F_7'$ top $m_3$ separately).} On class G/L the
Sklyanin tensor is an element of $\mathrm{Sym}^2(\fg)^{\fg}\otimes
\mathbb{Q}((z))$; its associator is a finite rational twist and
$\Phi_{17}\in\GRT^{\mathrm{fin}}$. On class M the correct additional
entry is $F_7'$ (the top $A_\infty$ operation), not the binary
Sklyanin face; its transport acquires MZV-coefficient corrections at
weight $\geq 3$ in the higher cochain tower
(Section~\ref{sec:f7-disambig}).

\smallskip
Uniqueness of $\Phi_{ij}$ follows from freeness of the
$\GRT$-action (Proposition~\ref{prop:grt-action-face}) after scalar
extension to a common coefficient field. The motivic and operadic
associators $\Phi_{18},\Phi_{19}$ are constructed in
Sections~\ref{sec:f8-motivic} and~\ref{sec:f9-willwacher} with the
coefficient-field scope of Convention~\ref{conv:grt-coefficient-field}.
\end{proof}

\begin{remark}[What this theorem replaces]
Vol I's seven-way chain of equivalences
$F_1\Leftrightarrow F_2\Leftrightarrow\cdots\Leftrightarrow F_7$
is absorbed, on the binary historical slice, as the statement that the
finite transports $\Phi_{1i}$ belong to $\GRT^{\mathrm{fin}}$. The
honest content, previously implicit,
is that the transports are associator-valued, not tautological
equalities; the chapter makes the associator visible and separates the
higher cochain faces from the binary residue.
\end{remark}

%%%%%%%%%%%%%%%%%%%%%%%%%%%%%%%%%%%%%%%%%%%%%%%%%%%%%%%%%%%%%%%%%%%%%%%%%%%%%
\section{F1\texorpdfstring{$\leftrightarrow$}{<->}F4 injection;
invisibility of MZV tail}
\label{sec:f1-f4}
%%%%%%%%%%%%%%%%%%%%%%%%%%%%%%%%%%%%%%%%%%%%%%%%%%%%%%%%%%%%%%%%%%%%%%%%%%%%%

The Drinfeld associator carries MZV data at all odd weights
$w\geq 3$. The residue $r(z)$ receives only the binary-degree
projection of that data.

\begin{theorem}[\label{thm:f1-f4-injection}F1$\leftrightarrow$F4 injection;
\ClaimStatusProvedHere]
Let $K$ be a characteristic-zero coefficient field containing the
coefficients of the two associators being compared. Let
$\Phi_{\KZ},\Phi_{\rat}\in\GRT(K(\!(\hbar)\!))$ be the
Knizhnik--Zamolodchikov and rational associator realisations after
scalar extension to $K$, and let $r(z)$ be the binary collision
residue. The map
$\mu\colon\GRT(K)\to
\cA^{!}\otimes\cA^{!}[\![z^{-1}]\!]$,
$\Phi\mapsto \Phi\cdot r(z)$, factors through the binary-degree
quotient $\GRT^{(2)}_K := \GRT(K)/\GRT^{(\geq 3)}(K)$, where
$\GRT^{(\geq 3)}$ is the pro-unipotent radical of associators
trivial at binary degree. Equivalently,
\[
 \Phi_{\KZ}\cdot r(z) \;=\; \Phi_{\rat}\cdot r(z);
\]
the MZV tail $\Phi_{\KZ}-\Phi_{\rat}$ is invisible to $r(z)$.
\end{theorem}

\begin{proof}
The collision residue $r(z)$ is homogeneous of bar-degree $2$, and
the associator action on bar-degree-$2$ elements factors through
the abelianisation of the universal enveloping filtration at
degree $\leq 2$. Write $\Phi = 1 + \sum_{w\geq 2}\hbar^w\phi_w$ with
$\phi_w\in\fgrt$ of $\fgrt$-weight $w$. The $\fgrt$-action on
$r(z)$ is
\[
 \Phi\cdot r(z) \;=\; r(z) + \sum_{w\geq 2}\hbar^w [\phi_w,r(z)]_{\mathrm{Ger}},
\]
where $[-,-]_{\mathrm{Ger}}$ is the Gerstenhaber bracket on
$\mathfrak{g}^{E_1}_{\cA}$. At bar-degree $2$ this bracket is
the classical Lie bracket on binary operations; $\phi_w$ for
$w\geq 3$ is supported on bar-degree $\geq 3$ and produces a
bar-degree $\geq 4$ correction. Projection to bar-degree $2$
annihilates all $w\geq 3$ contributions. Since $\Phi_{\KZ}-\Phi_{\rat}$
has zero weight-$2$ component (both associators are group-like with
identical quadratic part, the unique element of the pentagon axiom
at weight $2$), their action on $r(z)$ agrees. This proves the
factoring claim; the equivalent statement
$\Phi_{\KZ}\cdot r(z) = \Phi_{\rat}\cdot r(z)$ is the difference
vanishing.
\end{proof}

\begin{remark}[Content of Theorem~\ref{thm:f1-f4-injection}]
The associator carries an infinite tower of MZV coefficients
$\zeta(3),\zeta(5),\zeta(7),\dots$ (odd weights); the residue
$r(z)$ is a polynomial in $1/z$ with coefficients in
$\cA^{!}\otimes\cA^{!}$, and the MZV tower does not couple to the
binary-degree collision channel: the MZV invisibility is a
structural, not accidental, feature of the binary-degree truncation.
\end{remark}

%%%%%%%%%%%%%%%%%%%%%%%%%%%%%%%%%%%%%%%%%%%%%%%%%%%%%%%%%%%%%%%%%%%%%%%%%%%%%
\section{F7 disambiguation: Gaudin vs class-M top $A_\infty$ operation}
\label{sec:f7-disambig}
%%%%%%%%%%%%%%%%%%%%%%%%%%%%%%%%%%%%%%%%%%%%%%%%%%%%%%%%%%%%%%%%%%%%%%%%%%%%%

\begin{definition}[\label{def:f7-f7prime}$F_7$ vs $F_7'$]
\begin{enumerate}[(i)]
\item $F_7$ \emph{(Gaudin simple-pole)}: the canonical element
$t^{(2)}\in(\fg\otimes\fg)^{\fg}$ times the universal simple pole
$1/z$, giving the generator of the Gaudin Hamiltonian chain; defined
for all classes G/L/C/M.
\item $F_7'$ \emph{(class-M top $A_\infty$ $m_3$)}: the degree-$3$
component of the $A_\infty$ structure on $B(\cA)$ dual to the
Gerstenhaber algebra of the derived chiral centre; defined only
when the shadow tower has non-trivial $m_3$, i.e. on class M
chiral algebras with $r_{\max}\geq 3$.
\end{enumerate}
\end{definition}

\begin{proposition}[\label{prop:f7-split}F7 vs F7$'$ split;
\ClaimStatusProvedHere]
On class L (affine Kac--Moody), $F_7' = 0$ while
$F_7=t^{(2)}/z$ remains the binary Gaudin simple-pole face. Thus the
two symbols agree only after projecting to the higher bar-degree-three
lane where $F_7'$ lives: the projection of $F_7$ to that lane is zero.
On class M (Virasoro, $\Walg_N$), $F_7'$ is the chain-level obstruction
to constant-coproduct and is not a binary residue. Its simple-pole
shadow is $F_7$, while the genuine degree-three component is
\[
 \operatorname{gr}_3(F_7')
 \;=\;
 \sum_{w\geq 3} \hbar^w\,m_3^{(w)}
\]
in the bar-degree-$3$ component of $B(\cA)$, where each
$m_3^{(w)}$ is a weight-$w$ MZV-valued correction supplied by
$\phi_w\in\fgrt$ at bar-degree~$3$.
\end{proposition}

\begin{proof}
On class L the full $A_\infty$ structure on $B(V_k(\fg))$ truncates
at $m_2$ (Vol~I, prop:depth-gap-trichotomy, $d_{\mathrm{alg}}=1$).
Hence $m_3 = 0$ strictly, and $F_7' = 0$. The Gaudin element
$F_7 = t^{(2)}/z$ is the binary-degree-$2$ component and is
independent of $m_3$. On class M, $r_{\max} = \infty$ (Vol I,
class M fingerprint), so the ternary operation is a genuine higher
cochain datum. The displayed formula follows from the bar-degree
decomposition of $\Theta_{\cA}^{E_1}$: the bar-degree-$3$
Maurer--Cartan component is precisely the $A_\infty$ ternary
operation. Its binary simple-pole shadow is recorded by $F_7$, but
the higher operation itself is not equal to a binary residue. The MZV weight
structure of the $m_3^{(w)}$ is supplied by the $\fgrt$-action
through the $\phi_w$ generators.
\end{proof}

\begin{remark}[On the $\mathrm{F}_7$ ambiguity]
The ambiguity recorded earlier arose from the implicit
identification $F_7 = F_7'$. On class L the higher lane is zero, not
equal to the nonzero Gaudin residue. On class M the distinction is
substantive: $F_7$ is a simple-pole binary shadow and $F_7'$ is the
chain-level top operation carrying the MZV-weighted $m_3$ series. The
$\GRT$-torsor makes this difference visible as associator-dependent
data at bar-degree $3$.
\end{remark}

%%%%%%%%%%%%%%%%%%%%%%%%%%%%%%%%%%%%%%%%%%%%%%%%%%%%%%%%%%%%%%%%%%%%%%%%%%%%%
\section{Landscape-census chained-equality: Casimir rescaling
as \texorpdfstring{$\GRT^{\mathrm{fin}}$}{GRT-fin} gauge}
\label{sec:grtfin}
%%%%%%%%%%%%%%%%%%%%%%%%%%%%%%%%%%%%%%%%%%%%%%%%%%%%%%%%%%%%%%%%%%%%%%%%%%%%%

\begin{proposition}[\label{prop:grtfin-casimir}Casimir rescaling
and KZ level convention; tensor rescaling \ClaimStatusProvedHere,
KZ comparison \ClaimStatusConditional]
For affine Kac--Moody $V_k(\fg)$, the two common presentations of
$r(z)$,
\[
 r^{\mathrm{tr}}(z) \;=\; \frac{k\cdot\Omega_{\mathrm{tr}}}{z},
 \qquad
 r^{\KZ}(z) \;=\; \frac{\Omega}{(k+h^\vee)\,z},
\]
satisfy the tensor identity
\[
 \Omega_{\mathrm{tr}} \;=\; \frac{1}{2h^\vee}\,\Omega,
\]
where $\Omega_{\mathrm{tr}}$ is the trace-form Casimir and $\Omega$
is the Killing-normalised Casimir. The trace-form collision residue is
therefore
\[
 r^{\mathrm{tr}}(z)=\frac{k}{2h^\vee}\frac{\Omega}{z}.
\]
The comparison with $r^{\KZ}(z)$ is not a tensor identity in one level
coordinate; it is a KZ/Sugawara connection-convention comparison after
level normalisation. In particular the ``chained equality''
$k\cdot\Omega_{\mathrm{tr}}/z = \Omega/((k+h^\vee)z)$ is false as a
rational-function identity.
\end{proposition}

\begin{proof}
The trace form on $\fg$ is $\langle X,Y\rangle_{\mathrm{tr}}
= \operatorname{tr}_{\mathrm{def}}(XY)$ in the defining
representation; the Killing form is
$\langle X,Y\rangle_K = \operatorname{tr}(\operatorname{ad}X\cdot
\operatorname{ad}Y) = 2h^\vee\langle X,Y\rangle_{\mathrm{tr}}$
(standard Lie-algebraic identity). Hence
$\Omega = 2h^\vee\Omega_{\mathrm{tr}}$ (contracting with the
Killing form inverse rescales by the inverse factor), giving
$\Omega_{\mathrm{tr}} = \Omega/(2h^\vee)$. Substituting into
$r^{\mathrm{tr}}(z)$:
\[
 r^{\mathrm{tr}}(z) \;=\; \frac{k}{z}\cdot\frac{\Omega}{2h^\vee}
 \;=\; \frac{k}{2h^\vee}\cdot\frac{\Omega}{z}.
\]
Comparing with $r^{\KZ}(z) = \Omega/((k+h^\vee)z)$, the two
expressions agree only after changing both the Casimir normalisation
and the connection-level convention. The additive shift
$k\mapsto k+h^\vee$ is the Sugawara/KZ convention shift, not an
equality of tensors in the same coordinate. Both are valid coordinate
presentations in their respective Casimir bases and level conventions;
the bridge between them is conditional on naming that convention
change.
\end{proof}

\begin{remark}[Casimir rescaling as tensor identity]
The previous formulation wrote the tensor identity and the level
rescaling on the same line as a single equation, obscuring that two
distinct gauge transformations are in play. Proposition~\ref{prop:grtfin-casimir}
separates them: Casimir rescaling is the proved tensor identity, while
level rescaling is the Sugawara/KZ connection convention. The two
presentations are comparable only after both moves are named; they are
not the same Casimir in the same level convention.
\end{remark}

\begin{remark}[Compatibility with level-stripped $r$-matrix convention]
Proposition~\ref{prop:grtfin-casimir} preserves the level-stripped convention: in the
trace-form convention $r^{\mathrm{tr}}(z)|_{k=0}=0$; in the KZ
convention $r^{\KZ}(z)|_{k=0} = \Omega/(h^\vee z)\neq 0$ for
non-abelian $\fg$ (the Lie bracket persists abelian-limit). The
$\GRT^{\mathrm{fin}}$ transport carries one convention's boundary
behaviour to the other through the level rescaling, consistently
with the bridge identity
$k\Omega_{\mathrm{tr}} \leftrightarrow \Omega/(k+h^\vee)$ as a
connection-normalisation comparison at generic $k$.
\end{remark}

%%%%%%%%%%%%%%%%%%%%%%%%%%%%%%%%%%%%%%%%%%%%%%%%%%%%%%%%%%%%%%%%%%%%%%%%%%%%%
\section{Super-variant: odd generators and the $\mathbb{Z}/2$-factor}
\label{sec:super}
%%%%%%%%%%%%%%%%%%%%%%%%%%%%%%%%%%%%%%%%%%%%%%%%%%%%%%%%%%%%%%%%%%%%%%%%%%%%%

\begin{definition}[\label{def:grt-super}Super-$\GRT$]
For a chiral algebra $\cA$ with a $\mathbb{Z}/2$-graded generator
set $\{g_i\}$ with $|\mathrm{odd}|$ generators of odd parity, define
\[
 \Gsuper_K \;:=\; \GRT(K)\,\rtimes\,(\mathbb{Z}/2)^{|\mathrm{odd}|},
\]
where the semidirect factor records the parity-coordinate convention
on ordered half-edges in the collision-residue-to-Laplace comparison.
\end{definition}

\begin{proposition}[\label{prop:super-variant}Super-variant torsor;
\ClaimStatusConditional]
For a super-chiral algebra $\cA$, the face space
$\Fsuper(\cA)$ is a torsor over $\Gsuper_K$ after the parity-coordinate
convention is fixed. Heisenberg
(purely even, $p_{\max}=2$) and symplectic fermion (single odd
generator, $p_{\max}=1$) lie on distinct $\Gsuper_K$-orbits; the
separating invariant is the $\mathbb{Z}/2$-character of the
odd-odd OPE channel, not the algebra automorphism
$X\mapsto -X$ on both tensor legs.
\end{proposition}

\begin{proof}
The $\mathbb{Z}/2$-factor acts on the orientation line of the ordered
half-edge carrying an odd generator; this records the Koszul sign
appearing when the collision residue is compared with the Laplace
kernel. It is not the simultaneous generator sign
$X_i\mapsto(-1)^{|i|}X_i$ on both tensor legs, which would leave an
odd-odd tensor unchanged.
For Heisenberg ($|\mathrm{odd}|=0$) the odd sector is empty and
$r(z)$ is purely even; for symplectic fermion
($|\mathrm{odd}|=1$) the generator $\psi$ has OPE pole order $1$
(simple pole) and the collision residue is supported in the odd-odd
OPE channel of even total tensor parity. The two algebras have distinct collision-channel
characters under $\mathbb{Z}/2$, hence lie on distinct
$\Gsuper_K$-orbits. Transitivity of $\Gsuper_K$ on within-parity
classes follows from Proposition~\ref{prop:grt-action-face}
applied separately to each parity sector.
\end{proof}

\begin{remark}[Odd-generator collision residue sign]
For odd generators the collision residue
$r^{\coll}(z)$ differs from the Laplace-transform $r(z)$ by a
sign factor. In the super-variant formulation this sign is a
$\mathbb{Z}/2$-action on $\Face_K(\cA)$; the two residues are
realisations of the same $\Gsuper_K$-orbit in different
parity-sector coordinates. The Heisenberg-vs-symplectic-fermion
residue diagnostic is closed by Proposition~\ref{prop:super-variant}.
\end{remark}

%%%%%%%%%%%%%%%%%%%%%%%%%%%%%%%%%%%%%%%%%%%%%%%%%%%%%%%%%%%%%%%%%%%%%%%%%%%%%
\section{The motivic face \texorpdfstring{$F_8$}{F8}}
\label{sec:f8-motivic}
%%%%%%%%%%%%%%%%%%%%%%%%%%%%%%%%%%%%%%%%%%%%%%%%%%%%%%%%%%%%%%%%%%%%%%%%%%%%%

Brown's motivic programme (Brown 2012, arXiv:1102.1312) constructs
a coaction of the motivic Galois group $\mathcal{G}^{\mot}$ on the
ring $\mathcal{Z}^{\mot}$ of motivic multiple zeta values; this
coaction lifts to an action on the Drinfeld associator $\Phi_{\KZ}$
and hence, after scalar extension, on the higher cochain part of
$\Face_K(\cA)$.

\begin{definition}[\label{def:f8-motivic}Motivic face]
The \emph{motivic face} is
\[
 F_8 \;:=\; \Phi_{\mot}\cdot F_1,
\]
where $\Phi_{\mot}$ is the motivic associator viewed as an element of
$\GRT(\mathcal{Z}^{\mot})$, or after real period realisation as an
element of $\GRT(\mathbb{R})$. It is realised on the higher cochain
tower as
\[
 \Theta_{\geq3}^{\mot}
 \;=\;
 \Theta_{\geq3}^{\rat}
 \;+\;
 \sum_{w=3,5,7,\dots}\zeta^{\mot}(w)\,\Theta^{(w)}_{\geq3},
\]
while the binary residue $r(z)$ remains unchanged by
Theorem~\ref{thm:f1-f4-injection}.
\end{definition}

\begin{theorem}[\label{thm:motivic-face-F8}Motivic face is a $\GRT$-orbit;
\ClaimStatusConditional]
The motivic face $F_8$ is a well-defined higher cochain point of
$\Face_{\mathcal Z^{\mot}}(\cA)$, and after real period realisation a
point of $\Face_{\mathbb R}(\cA)$. Its binary residue agrees with
$F_1$; its non-trivial content begins in bar degree at least three.
\end{theorem}

\begin{proof}
The motivic coaction of Brown 2012 (Thm.~1.1) gives a Hopf
algebra coaction
$\Delta^{\mot}\colon\mathcal{Z}^{\mot}\to
\mathcal{Z}^{\mot}\otimes(\mathcal{Z}^{\mot}/\zeta(2))$, inducing
a $\mathcal{G}^{\mot}$-module structure on $\mathcal{Z}^{\mot}$.
The KZ associator $\Phi_{\KZ}$ has motivic lift $\Phi_{\KZ}^{\mot}
\in\GRT(\mathcal{Z}^{\mot})$ (Brown 2012, \S3); composition with
the canonical realisation $\mathcal{Z}^{\mot}\to\mathbb{R}$
supplies a real point of $\GRT(\mathbb R)$, not a rational point of
$\GRT(\mathbb Q)$. Applying
Proposition~\ref{prop:grt-action-face} over
$K=\mathcal{Z}^{\mot}$ or $K=\mathbb R$, the element
$\Phi_{\mot}\cdot F_1$ is a well-defined higher cochain face. By
Theorem~\ref{thm:f1-f4-injection}, the binary-degree component of
$F_8$ is identical to that of $F_1$; the non-triviality of $F_8$ lies
entirely at bar-degree $\geq 3$, carrying the MZV-weighted corrections
$\zeta^{\mot}(w)\,\Theta^{(w)}_{\geq3}$.
\end{proof}

\begin{remark}[Odd-weight generation of $\fgrt$]
The conjectural structure of $\fgrt$ (Drinfeld, Kassel--Turaev)
says $\fgrt \simeq \mathrm{Lie}(\sigma_3,\sigma_5,\sigma_7,\dots)$
with odd-weight generators; the motivic lift by Brown is the
universal construction realising this as a
$\mathcal{G}^{\mot}$-module. Each generator $\sigma_{2k+1}$
($k\geq 1$) supplies an independent MZV tail at weight $2k+1$,
which in Definition~\ref{def:f8-motivic} enters as the
coefficient of $\zeta^{\mot}(2k+1)$ in the higher cochain tower
$\Theta_{\geq3}^{\mot}$.
\end{remark}

%%%%%%%%%%%%%%%%%%%%%%%%%%%%%%%%%%%%%%%%%%%%%%%%%%%%%%%%%%%%%%%%%%%%%%%%%%%%%
\section{The operadic face \texorpdfstring{$F_9$}{F9}}
\label{sec:f9-willwacher}
%%%%%%%%%%%%%%%%%%%%%%%%%%%%%%%%%%%%%%%%%%%%%%%%%%%%%%%%%%%%%%%%%%%%%%%%%%%%%

Willwacher's theorem (Willwacher 2015, arXiv:1009.1654) identifies
$H^0$ of Kontsevich's graph complex $\GC_2$ with $\fgrt$; this
gives the $\fgrt$-action on the deformation complex of the $E_2$
operad, and via Tamarkin chain, on $B(\cA)$.

\begin{definition}[\label{def:f9-willwacher}Operadic face]
The \emph{operadic face} is
\[
 F_9 \;:=\; \Phi_{\mathrm{Wil}}\cdot F_1,
\]
where $\Phi_{\mathrm{Wil}}\in\GRT(\mathbb{Q})$ is the associator
whose Lie-algebra element $\log\Phi_{\mathrm{Wil}}\in\fgrt$ is a
Willwacher graph cocycle representative (a bivalent-graph cocycle
in $\GC_2$ of odd loop-order).
\end{definition}

\begin{theorem}[\label{thm:willwacher-face-F9}Willwacher face is a
$\GRT$-orbit; \ClaimStatusConditional]
After fixing the rational graph-complex model and the Tamarkin chain,
the Willwacher face $F_9$ is a well-defined higher cochain point of
$\Face_{\mathbb Q}(\cA)$.
The Tamarkin chain
\[
 \GC_2 \;\xrightarrow{\ \simeq\ }\; \Def(E_2)
 \;\xrightarrow{\ \simeq\ }\; \Def\bigl(B(-)\bigr)
\]
realises $\log\Phi_{\mathrm{Wil}}\in H^0(\GC_2) = \fgrt$ as an
infinitesimal deformation of the bar differential on $B(\cA)$;
exponentiating gives the $\GRT(\mathbb Q)$-element
$\Phi_{\mathrm{Wil}}$ which acts on $F_1$ to produce $F_9$.
\end{theorem}

\begin{proof}
Willwacher 2015, Thm.~1.1, gives the isomorphism
$H^0(\GC_2)\simeq\fgrt$. The Tamarkin chain (Tamarkin 1998,
reformulated in Willwacher 2015 \S9) is a zigzag of
quasi-isomorphisms of $L_\infty$-algebras connecting $\GC_2$
to the deformation complex $\Def(E_2)$ of the $E_2$-operad, and
the canonical formality functor
$\Def(E_2)\xrightarrow{\sim}\Def(B(-))$ transports $\fgrt$-action
to the deformation complex of the bar construction. A Willwacher
graph cocycle $\gamma\in\GC_2$ of loop order $2k+1$ (odd-loop
cocycles span $\fgrt$ at weight $2k+1$) defines an infinitesimal
deformation of the bar differential on $B(\cA)$ along the
Tamarkin chain; exponentiating to the group-level gives
$\Phi_{\mathrm{Wil}} := \exp(\gamma)\in\GRT(\mathbb{Q})$. Applying
Proposition~\ref{prop:grt-action-face}, the element
$\Phi_{\mathrm{Wil}}\cdot F_1$ is a well-defined face, which we
denote $F_9$. The transport is independent of the choice of
cocycle representative up to exact terms, since the isomorphism
$H^0(\GC_2)\simeq\fgrt$ is canonical.
\end{proof}

\begin{remark}[Tamarkin chain as motivic--operadic bridge]
The Tamarkin chain converts a graph-complex cocycle (operadic
data) into an associator (algebraic data). It is the structural
bridge between $F_8$ (motivic, geometric) and $F_9$ (operadic,
combinatorial). Both faces enter the torsor at the same
$\fgrt$-generator level and are identified up to associator gauge.
\end{remark}

%%%%%%%%%%%%%%%%%%%%%%%%%%%%%%%%%%%%%%%%%%%%%%%%%%%%%%%%%%%%%%%%%%%%%%%%%%%%%
\section{Duality corollary: $F_8$ and $F_9$ as geometric vs operadic
incarnations of one $\GRT$-orbit}
\label{sec:f8-f9-dual}
%%%%%%%%%%%%%%%%%%%%%%%%%%%%%%%%%%%%%%%%%%%%%%%%%%%%%%%%%%%%%%%%%%%%%%%%%%%%%

\begin{corollary}[\label{cor:f8-f9-dual}F8--F9 duality;
\ClaimStatusConditional]
After scalar extension to a coefficient field containing both the
motivic-period realisation of $\Phi_{\mot}$ and the rational
Willwacher representative $\Phi_{\mathrm{Wil}}$, the faces $F_8$ and
$F_9$ lie on the same $\GRT$-orbit. The connecting associator is
\[
 \Phi_{89} = \Phi_{\mathrm{Wil}}\cdot\Phi_{\mot}^{-1}
\]
in that common scalar extension. This is not a claim that
$\Phi_{\mot}$ is a rational point of $\GRT(\mathbb Q)$.
\end{corollary}

\begin{proof}
By Theorems~\ref{thm:motivic-face-F8} and~\ref{thm:willwacher-face-F9},
$F_8 = \Phi_{\mot}\cdot F_1$ and $F_9 = \Phi_{\mathrm{Wil}}\cdot F_1$.
Hence
$F_9 = (\Phi_{\mathrm{Wil}}\Phi_{\mot}^{-1})\cdot F_8$, and the
connecting associator is
$\Phi_{89} := \Phi_{\mathrm{Wil}}\Phi_{\mot}^{-1}$ in the common
coefficient field.
The Tamarkin chain sends Brown's motivic associator to a graph
cocycle representative of the same $\fgrt$-class (Willwacher 2015,
\S9; Deligne--Terasoma lift), realising the identification at the
infinitesimal level. The orbit-equivalence follows from the free
and transitive $\GRT$-action on $\Face_K(\cA)$.
\end{proof}

\begin{remark}[Geometric vs operadic]
$F_8$ is ``geometric'' in the sense that its construction uses
configuration spaces of points on $\mathbb{P}^1$ (iterated-integral
realisation of MZVs). $F_9$ is ``operadic'' in the sense that its
construction uses graph complexes (combinatorial realisation of
$\fgrt$). Corollary~\ref{cor:f8-f9-dual} records that the two
constructions are two coordinates on the same torsor orbit; the
geometry/operad duality is encoded in the Tamarkin chain.
\end{remark}

%%%%%%%%%%%%%%%%%%%%%%%%%%%%%%%%%%%%%%%%%%%%%%%%%%%%%%%%%%%%%%%%%%%%%%%%%%%%%
\section{Synopsis and external references}
\label{sec:grt-synopsis}
%%%%%%%%%%%%%%%%%%%%%%%%%%%%%%%%%%%%%%%%%%%%%%%%%%%%%%%%%%%%%%%%%%%%%%%%%%%%%

The chapter upgrades the Vol~I seven-face programme to a
$\GRT(K)$-torsor $\Face_K(\cA)$ after a coefficient field and formality
base point are fixed, with historical rational finite-gauge coordinates
$F_1,\dots,F_7$, two higher cochain additions $F_8$ (Brown motivic
period) and $F_9$ (Willwacher operadic), and an explicit super-variant
$\Fsuper$. The master theorem is
Theorem~\ref{thm:seven-faces-as-GRT-torsor}, conditional on the
coefficient-field and concrete-face scopes of
Convention~\ref{conv:grt-coefficient-field}. The three closed
confusions are:
Casimir rescaling vs chained equality
(Proposition~\ref{prop:grtfin-casimir}),
F1$\leftrightarrow$F4 MZV-invisibility
(Theorem~\ref{thm:f1-f4-injection}),
F7 vs F7$'$ split (Proposition~\ref{prop:f7-split}),
and the coefficient-field split for motivic versus operadic faces
(Convention~\ref{conv:grt-coefficient-field}).

External anchors:
\begin{itemize}
\item Drinfeld 1991 ``On quasitriangular quasi-Hopf algebras'',
Leningrad Math.~J.~2: the definition and basic properties of
$\GRT$.
\item Brown 2012, arXiv:1102.1312, ``Mixed Tate motives over
$\mathbb{Z}$'': the motivic coaction on MZVs and the motivic
associator.
\item Willwacher 2015, arXiv:1009.1654, ``M.~Kontsevich's graph
complex and the Grothendieck--Teichm\"uller Lie algebra'':
$H^0(\GC_2)\simeq\fgrt$ and the action on $\Def(E_2)$.
\item Tamarkin 1998, ``Another proof of M.~Kontsevich formality
theorem'': the Tamarkin chain
$\GC_2\to\Def(E_2)\to\Def(B(-))$.
\item Loday--Vallette, \emph{Algebraic Operads}, Ch.~11--12: the
bar-cobar framework and twisting morphisms.
\end{itemize}

These three anchors are the independent-verification sources of
record for Theorems~\ref{thm:seven-faces-as-GRT-torsor},
\ref{thm:motivic-face-F8}, \ref{thm:willwacher-face-F9},
and~\ref{cor:f8-f9-dual}. Each is disjoint from the derivation
route used in the chapter proofs: Drinfeld 1991 supplies the
$\GRT$ group-theoretic backbone, Brown 2012 supplies the motivic
associator independently of graph-complex machinery, and
Willwacher 2015 supplies the operadic realisation independently of
motivic machinery. The decorator declarations attached to the
test infrastructure
(\texttt{compute/tests/test\_grt\_parametrized\_seven\_faces.py})
encode these disjointness relations formally.

%%%%%%%%%%%%%%%%%%%%%%%%%%%%%%%%%%%%%%%%%%%%%%%%%%%%%%%%%%%%%%%%%%%%%%%%%%%%%
\section{Standing qualifiers: associator, scope, slab-bimodule}
\label{sec:grt-qualifiers}
%%%%%%%%%%%%%%%%%%%%%%%%%%%%%%%%%%%%%%%%%%%%%%%%%%%%%%%%%%%%%%%%%%%%%%%%%%%%%

The $\GRT$-torsor master Theorem~\ref{thm:seven-faces-as-GRT-torsor}
inherits three standing scope qualifiers. Recording them here pins
the precise content of the torsor statement and connects the chapter
to the climax theorems.

\begin{remark}[Cohomology vs.\ cochain associator action]
\label{rem:grt-cohomology-vs-cochain}
The $\GRT(K)$-action on $\Face_K(\cA)$ of
Proposition~\ref{prop:grt-action-face} is a strict groupoid action
on the associator-coordinate component after the coefficient field is
fixed; this is the content of Tamarkin's theorem combined with
Kontsevich--Tamarkin formality.
At the level of cohomology classes in
$H^{*}(\cZ^{\mathrm{der}}_{\mathrm{ch}}(\cA))$, however, the
$\GRT$-action factors through the abelianisation
$\GRT^{\mathrm{ab}} = \fgrt^{\mathrm{ab}}$, which is the
$\bQ$-vector space of Drinfeld weight-$\geq 2$ generators modulo
Lie commutators; the binary-degree projection
of Theorem~\ref{thm:f1-f4-injection} kills the entire
$\fgrt$-bracket-content of $\Phi$. Hence:
\begin{itemize}
\item[(a)] \emph{Cohomological face} $r(z)$ \emph{is associator-independent.}
 Two associators $\Phi_1, \Phi_2 \in \GRT(K)$ produce the
 same $r(z)\in H^{*}(\cZ^{\mathrm{der}}_{\mathrm{ch}}(\cA))^{\otimes 2}$
 binary-degree class.
\item[(b)] \emph{Cochain-level face} $r(z)|^{\Phi}$
 \emph{depends on $\Phi$.} The choice of $\Phi$ fixes a strict
 representative; the transport between $r(z)|^{\Phi_1}$ and
 $r(z)|^{\Phi_2}$ is an explicit cochain-level associator transport
 along the Tamarkin chain (Section~\ref{sec:f9-willwacher}).
\end{itemize}
The cohomological reading is the canonical one for cross-volume
identifications (Vol I bar invariants, Vol III $\Phi$ functor); the
cochain-level reading is the technical realisation needed to compute
explicit higher Massey products (e.g.\ the F7$'$ class-M top
$A_\infty$ operation of Proposition~\ref{prop:f7-split}, where MZV
weights $\geq 3$ enter as cochain-level associator data).
\end{remark}

\begin{remark}[Abstract Koszul-locus scope of the torsor theorem]
\label{rem:grt-abstract-vs-concrete}
Theorem~\ref{thm:seven-faces-as-GRT-torsor} is proved abstractly only
after fixing $K$ and a formality base point on the chirally Koszul
locus for any
$\SCchtop$-algebra in the standard 3d HT landscape: the
$\GRT(K)$-torsor structure is a consequence of Tamarkin's formality
and is independent of the specific algebra. The
\emph{concrete} construction of all face transports
$\Phi_{1i}$ as explicit rational finite-gauge associators is verified
only for the binary historical coordinates in the $\fsl_2$-Yangian and
affine KM at non-critical level; F8 and F9 require the scalar
extensions of Convention~\ref{conv:grt-coefficient-field};
the elliptic case is partial (KZB associator constructed at the
classical level, full quantum-$R$-matrix
quasi-triangularity at finite $\hbar$ open); the toroidal case is
absent at the present scope. See
Remark~\ref{rem:concrete-vs-abstract-scope} of the
unified-chiral-quantum-group chapter and
Remark~\ref{rem:dnp-master-concrete-vs-abstract} of the
DNP-identification-master chapter.
\end{remark}

\begin{remark}[Slab is a bimodule: $r(z)$ as diagonal-bimodule data]
\label{rem:grt-slab-bimodule-scope}
The face space $\Face_K(\cA)$ enters 3d HT physics as the space of
chain-level presentations of the $r$-matrix datum on the
\emph{diagonal} of the Dimofte $D|T|N$ slab bimodule
$(\cA_{\mathrm{L}} = \cA_{\mathrm{R}} = \cA)$. The off-diagonal slab
carries an $r$-matrix in
$\cA_{\mathrm{L}}^{!}\otimes\cA_{\mathrm{R}}^{!}[\![z^{-1}]\!]$ that
mediates the categorified bimodule structure; the corresponding
face space $\Face(\cA_{\mathrm{L}}, \cA_{\mathrm{R}})$ is a torsor
over a generalised $\GRT$-groupoid. The diagonal restriction
recovers the present chapter's $\GRT(K)$-torsor on
$\Face_K(\cA, \cA) = \Face_K(\cA)$ after scalar extension.
\end{remark}

\begin{remark}[Cross-references to the climax]
\label{rem:grt-climax-anchors}
The torsor master Theorem~\ref{thm:seven-faces-as-GRT-torsor}
threads three reconstitution anchors:
\begin{itemize}
\item[(a)] \emph{Universal holography master.} Each face
 $F_i \in \Face_K(\cA)$ is a coordinate presentation of the
 boundary-to-bulk obstruction $r(z)$ in the universal-holography
 functor $\Phi_{\mathrm{hol}}\colon \cA\mapsto T_{\cA}$
 (Theorem~\ref{thm:universal-holography-master}). The five rungs
 (Corollaries~\ref{cor:rung-heisenberg}, \ref{cor:rung-affine-km},
 \ref{cor:rung-virasoro}, \ref{cor:rung-w-N},
 \ref{cor:rung-w-infty}) each produce a $\GRT(K)$-torsor
 of seven-face presentations via the present theorem; the $\fgrt$
 generators $\sigma_3, \sigma_5, \dots$ specialise to the higher-spin
 stress tensors of the $E_{n}$-topologisation ladder
 (Theorem~\ref{thm:e-infinity-topologization-ladder}).
\item[(b)] \emph{Universal conductor.} The Volume~I universal-conductor
 $K = c + c^{!}$ is the $\GRT^{\mathrm{fin}}$-invariant scalar
 pinning the absolute Casimir normalisation of $F_1$. The conductor
 formula expresses the ghost charge as the trace of a specific
 $\fgrt^{\mathrm{ab}}$-cocycle representative on the binary-degree
 bar component, recovering the present chapter's
 $\GRT^{\mathrm{fin}}$-classification of $F_1, \dots, F_7$ as a
 refinement of the conductor's Casimir-rescaling axis.
\item[(c)] \emph{Volume~III $\Phi$ functor.} The Volume~III
 CY-to-chiral functor $\Phi$ pulls back the present $\GRT$-torsor
 structure to $\mathrm{Coh}(X)$: each CY target carries a seven-face
 torsor associated with $\Phi(X)$, with explicit identifications at
 K3 (via the Mukai-orthogonal Yangian
 $Y_{\hbar}(\mathfrak{so}(4,20))$, with conditional
 $Y_{\hbar}(\mathfrak{so}(4|20))$ Hodge-parity refinement), $\bC^{3}$
 (via Schiffmann--Vasserot CoHA), and toric CY${}_{3}$ (via
 Maulik--Okounkov $r$-matrix). The
 Tamarkin chain of Section~\ref{sec:f9-willwacher} is the operadic
 bridge; the motivic face $F_8$ specialises in Volume~III to the
 Brown motivic associator on the moduli of CY-period configurations.
 Scope: $n=2$ for $d\le 2$, $n=1$ for $d\ge 3$.
\end{itemize}
\end{remark}

%%%%%%%%%%%%%%%%%%%%%%%%%%%%%%%%%%%%%%%%%%%%%%%%%%%%%%%%%%%%%%%%%%%%%%%%%%%%%
%%% End of chapter.
%%%%%%%%%%%%%%%%%%%%%%%%%%%%%%%%%%%%%%%%%%%%%%%%%%%%%%%%%%%%%%%%%%%%%%%%%%%%%
